%! Unit 2: Solving Linear Equations Ax = b — Summative Assessment — Practice Test --- Study Copy (KEY)
\documentclass[10pt]{article}
\usepackage{linalg-article}
\usepackage{linalg-key}   % pulls in -boxes; do NOT also load -boxes
\newcommand{\bb}{\mathbf{b}}
\newcommand{\xx}{\mathbf{x}}
\newcommand{\cc}{\mathbf{c}}

% Part-divider strip
\newcommand{\parthead}[1]{%
  \vspace{6pt}\noindent
  \begin{headlinebox}{burgundy}{\color{white}\bfseries #1}\end{headlinebox}%
  \vspace{4pt}}

\begin{document}

\pageheader{Unit 2: Solving Linear Equations Ax = b}{Practice Test --- Study Copy --- Answer Key}
\namedateperiod

\vspace{0.10in}
\begin{remindbox}
\small
\textbf{This is a practice test.} It mirrors the real Unit 2 test in length, format,
and the ideas it covers, but the numbers and contexts differ. Work every problem
with no notes, then check your answers against the key.
\end{remindbox}

% ======================================================================
\parthead{Part A --- Vocabulary (8 pts)}
Write the letter of the best description next to each term.

\vspace{4pt}
\noindent
\begin{minipage}[t]{0.44\textwidth}
\begin{enumerate}[leftmargin=2.2em, itemsep=7pt, label=\arabic*.]
  \item Pivot \ans{B}
  \item Multiplier \ans{D}
  \item Upper triangular ($U$) \ans{G}
  \item Inverse matrix \ans{H}
  \item Singular matrix \ans{C}
  \item $LU$ factorization \ans{A}
  \item Permutation matrix \ans{E}
  \item Transpose \ans{F}
\end{enumerate}
\end{minipage}%
\hfill
\begin{minipage}[t]{0.53\textwidth}
\small
\begin{description}[leftmargin=1.4em, itemsep=3pt]
  \item[(A)] $A=LU$: a lower-triangular $L$ times an upper-triangular $U$.
  \item[(B)] The first nonzero entry in a row, used to clear the entries below it.
  \item[(C)] A matrix with no inverse --- its determinant $ad-bc$ is $0$.
  \item[(D)] $\ell=(\text{entry to clear})\div(\text{pivot})$; how many pivot rows you subtract.
  \item[(E)] The identity with two rows swapped; it records a row exchange.
  \item[(F)] The matrix $A^{\mathsf T}$ formed by flipping $A$ across its diagonal (rows $\leftrightarrow$ columns).
  \item[(G)] A matrix with all zeros below the main diagonal.
  \item[(H)] $A^{-1}$, the matrix that undoes $A$: $A^{-1}A=I$.
\end{description}
\end{minipage}

% ======================================================================
\parthead{Part B --- Multiple Choice (12 pts)}
Circle the best answer.

\begin{enumerate}[leftmargin=1.9em, itemsep=9pt]
  \item To clear the entry below a pivot $p$, you subtract $\ell$ times the pivot row, where $\ell$ is
    \hfill (A) $p\times(\text{entry})$ \quad \textbf{(B) $(\text{entry})\div p$} \ans{$\leftarrow$} \quad
    (C) $p\div(\text{entry})$ \quad (D) $p-(\text{entry})$
  \item In the factorization $A=LU$, the matrix $L$ holds
    \hfill (A) the pivots \quad (B) the solution $\xx$ \\
    \phantom{x}\hfill \textbf{(C) the multipliers, with $1$'s on the diagonal} \ans{$\leftarrow$} \quad (D) all zeros
  \item A $2\times2$ matrix $A=\begin{bsmallmatrix}a&b\\c&d\end{bsmallmatrix}$ is \textbf{singular} (no inverse) exactly when
    \hfill \textbf{(A) $ad-bc=0$} \ans{$\leftarrow$} \quad (B) $ad-bc=1$ \quad (C) $a=d$ \quad (D) it is upper triangular
  \item If $P$ is a permutation matrix, the product $PA$
    \hfill (A) scales the rows of $A$ \quad \textbf{(B) swaps rows of $A$} \ans{$\leftarrow$} \quad (C) transposes $A$ \quad (D) inverts $A$
  \item For any two matrices that can be multiplied, $(AB)^{\mathsf T}=$
    \hfill (A) $A^{\mathsf T}B^{\mathsf T}$ \quad (B) $AB$ \quad (C) $BA$ \quad \textbf{(D) $B^{\mathsf T}A^{\mathsf T}$} \ans{$\leftarrow$}
  \item Two equations in two unknowns have \textbf{infinitely many} solutions when the two lines
    \hfill (A) cross once \quad (B) are parallel but different \quad \textbf{(C) are the same line} \ans{$\leftarrow$} \quad (D) are perpendicular
\end{enumerate}

\begin{teachernote}
Item 2: $L$ records the elimination multipliers (with unit diagonal); $U$ holds the pivots. Item 5:
socks-and-shoes order reverses --- $(AB)^{\mathsf T}=B^{\mathsf T}A^{\mathsf T}$, not $A^{\mathsf T}B^{\mathsf T}$.
Item 6: same line $=$ one equation is a multiple of the other (dependent) $\Rightarrow$ infinitely many;
parallel-but-different $\Rightarrow$ none.
\end{teachernote}

% ======================================================================
\parthead{Part C --- Short Answer \& Computation (30 pts)}
Show your work.

\begin{enumerate}[leftmargin=1.9em, itemsep=12pt]
  \item Solve the system by \textbf{elimination}, then back-substitute and check. State the pivot and the
        multiplier you used.
        \[ x+2y=4, \qquad 3x+4y=10. \]
        \ans{Pivot $=1$ (the $x$-coefficient of eq.~1); multiplier $\ell=3\div1=3$. Eq.~2 $-\,3\times$ eq.~1: $(4-6)y=10-12$, so $-2y=-2$, $y=1$. Back-substitute: $x+2(1)=4$, $x=2$. Solution $(x,y)=(2,1)$. Check: $3(2)+4(1)=10$ \checkmark.}
  \item Row-reduce $A=\begin{bsmallmatrix}1&2&1\\2&5&3\\4&10&8\end{bsmallmatrix}$ to upper-triangular $U$.
        List the three multipliers and the pivots.
        \ans{$R_2-2R_1$ ($\ell_{21}=2$) and $R_3-4R_1$ ($\ell_{31}=4$) give $\begin{bsmallmatrix}1&2&1\\0&1&1\\0&2&4\end{bsmallmatrix}$; then $R_3-2R_2$ ($\ell_{32}=2$) gives $U=\begin{bsmallmatrix}1&2&1\\0&1&1\\0&0&2\end{bsmallmatrix}$. Multipliers $\ell_{21}=2,\ \ell_{31}=4,\ \ell_{32}=2$; pivots $1,\,1,\,2$.}
  \item Let $A=\begin{bsmallmatrix}2&1\\3&2\end{bsmallmatrix}$. Use the $ad-bc$ formula to find $A^{-1}$,
        then check that $A^{-1}A=I$.
        \ans{$ad-bc=2\cdot2-1\cdot3=1$, so $A^{-1}=\dfrac{1}{1}\begin{bsmallmatrix}2&-1\\-3&2\end{bsmallmatrix}=\begin{bsmallmatrix}2&-1\\-3&2\end{bsmallmatrix}$. Check: $\begin{bsmallmatrix}2&-1\\-3&2\end{bsmallmatrix}\begin{bsmallmatrix}2&1\\3&2\end{bsmallmatrix}=\begin{bsmallmatrix}1&0\\0&1\end{bsmallmatrix}$ \checkmark.}
  \item Using the inverse from the previous problem, solve $A\xx=\bb$ for
        $\bb=\begin{bsmallmatrix}3\\4\end{bsmallmatrix}$ (compute $\xx=A^{-1}\bb$).
        \ans{$\xx=A^{-1}\bb=\begin{bsmallmatrix}2&-1\\-3&2\end{bsmallmatrix}\begin{bsmallmatrix}3\\4\end{bsmallmatrix}=\begin{bsmallmatrix}6-4\\-9+8\end{bsmallmatrix}=\begin{bsmallmatrix}2\\-1\end{bsmallmatrix}$. Check: $\begin{bsmallmatrix}2&1\\3&2\end{bsmallmatrix}\begin{bsmallmatrix}2\\-1\end{bsmallmatrix}=\begin{bsmallmatrix}3\\4\end{bsmallmatrix}$ \checkmark.}
  \item Factor $A=\begin{bsmallmatrix}2&4\\6&7\end{bsmallmatrix}$ as $A=LU$. Write $L$ and $U$, then verify
        $LU=A$.
        \ans{Pivot $2$, multiplier $\ell=6\div2=3$. $R_2-3R_1$: $\begin{bsmallmatrix}0&7-12\end{bsmallmatrix}=\begin{bsmallmatrix}0&-5\end{bsmallmatrix}$, so $U=\begin{bsmallmatrix}2&4\\0&-5\end{bsmallmatrix}$, $L=\begin{bsmallmatrix}1&0\\3&1\end{bsmallmatrix}$. Check $LU=\begin{bsmallmatrix}1&0\\3&1\end{bsmallmatrix}\begin{bsmallmatrix}2&4\\0&-5\end{bsmallmatrix}=\begin{bsmallmatrix}2&4\\6&7\end{bsmallmatrix}=A$ \checkmark.}
  \item Using the $L$ and $U$ from the previous problem, solve $A\xx=\bb$ for
        $\bb=\begin{bsmallmatrix}2\\1\end{bsmallmatrix}$ in \textbf{two sweeps}: first $L\cc=\bb$
        (forward), then $U\xx=\cc$ (back).
        \ans{$L\cc=\bb$: $c_1=2$; $3c_1+c_2=1\Rightarrow c_2=1-6=-5$, so $\cc=\begin{bsmallmatrix}2\\-5\end{bsmallmatrix}$. $U\xx=\cc$: $-5x_2=-5\Rightarrow x_2=1$; $2x_1+4=2\Rightarrow x_1=-1$, so $\xx=\begin{bsmallmatrix}-1\\1\end{bsmallmatrix}$. Check: $\begin{bsmallmatrix}2&4\\6&7\end{bsmallmatrix}\begin{bsmallmatrix}-1\\1\end{bsmallmatrix}=\begin{bsmallmatrix}2\\1\end{bsmallmatrix}$ \checkmark.}
  \item \textbf{(a)} The matrix $A=\begin{bsmallmatrix}0&2\\1&3\end{bsmallmatrix}$ has a $0$ in the top-left
        pivot spot. Write the permutation matrix $P$ that swaps the two rows, and give $PA$.
        \textbf{(b)} Write the transpose of $B=\begin{bsmallmatrix}1&2\\2&5\end{bsmallmatrix}$ and state
        whether $B$ is symmetric.
        \ans{(a) $P=\begin{bsmallmatrix}0&1\\1&0\end{bsmallmatrix}$, and $PA=\begin{bsmallmatrix}1&3\\0&2\end{bsmallmatrix}$ (now a nonzero pivot sits in the corner). (b) $B^{\mathsf T}=\begin{bsmallmatrix}1&2\\2&5\end{bsmallmatrix}=B$, so $B$ \emph{is} symmetric.}
\end{enumerate}

% ======================================================================
\parthead{Part D --- Extended Response (10 pts)}
Explain your reasoning in full sentences.

\begin{enumerate}[leftmargin=1.9em, itemsep=12pt]
  \item Let $A=\begin{bsmallmatrix}1&2\\2&4\end{bsmallmatrix}$.
        \textbf{(a)} Compute $ad-bc$. Does $A$ have an inverse? What does that say about solving $A\xx=\bb$?
        \textbf{(b)} How many solutions does $A\xx=\begin{bsmallmatrix}3\\6\end{bsmallmatrix}$ have?
        \textbf{(c)} How many does $A\xx=\begin{bsmallmatrix}1\\1\end{bsmallmatrix}$ have? Justify (b) and (c)
        using elimination or the row/line picture.
        \ans{(a) $ad-bc=1\cdot4-2\cdot2=0$, so $A$ is \emph{singular} --- no inverse. Elimination hits a zero pivot ($R_2-2R_1$ wipes out row~2), so $A\xx=\bb$ has \emph{either} no solution \emph{or} infinitely many, never a unique one. (b) Apply $R_2-2R_1$ to $\bb$: $6-2(3)=0$, giving $0=0$ --- a true statement, so \emph{infinitely many} solutions (the two equations are the same line). (c) $R_2-2R_1$ on $\begin{bsmallmatrix}1\\1\end{bsmallmatrix}$: $1-2(1)=-1$, giving $0=-1$ --- impossible, so \emph{no solution} (parallel but different lines).}
  \item Let $A=\begin{bsmallmatrix}1&2\\3&1\end{bsmallmatrix}$.
        \textbf{(a)} Compute $A^{\mathsf T}A$ and check that it is symmetric.
        \textbf{(b)} Explain why $A^{\mathsf T}A$ is \emph{always} symmetric for any matrix $A$, using the
        reversal rule $(AB)^{\mathsf T}=B^{\mathsf T}A^{\mathsf T}$.
        \ans{(a) $A^{\mathsf T}=\begin{bsmallmatrix}1&3\\2&1\end{bsmallmatrix}$, so $A^{\mathsf T}A=\begin{bsmallmatrix}1&3\\2&1\end{bsmallmatrix}\begin{bsmallmatrix}1&2\\3&1\end{bsmallmatrix}=\begin{bsmallmatrix}10&5\\5&5\end{bsmallmatrix}$ --- the off-diagonal entries match, so it is symmetric. (b) Take the transpose of $A^{\mathsf T}A$ and apply the reversal rule: $(A^{\mathsf T}A)^{\mathsf T}=A^{\mathsf T}(A^{\mathsf T})^{\mathsf T}=A^{\mathsf T}A$, since $(A^{\mathsf T})^{\mathsf T}=A$. A matrix equal to its own transpose is symmetric, so $A^{\mathsf T}A$ is always symmetric.}
\end{enumerate}

\begin{teachernote}
\textbf{Scoring Part D (5 pts each).} D1: 1 pt (a) $ad-bc=0\Rightarrow$ singular/no inverse; 2 pts (b)
infinitely many with the $0=0$ justification; 2 pts (c) none with the $0=-1$ justification. Full credit
needs the \emph{reason} (same line vs.\ parallel), not just a count. D2: 2 pts (a) correct
$A^{\mathsf T}A=\begin{bsmallmatrix}10&5\\5&5\end{bsmallmatrix}$ and noting equal off-diagonals; 3 pts (b)
the general argument $(A^{\mathsf T}A)^{\mathsf T}=A^{\mathsf T}A$ via the reversal rule. Watch for students
writing $(AB)^{\mathsf T}=A^{\mathsf T}B^{\mathsf T}$ (wrong order).
\end{teachernote}

\end{document}
